% The paper's brickwork graph, with the general MBQC points revealed below.
\begin{frame}{Background: Measurement-Based Quantum Computation}
  \centering
  \setbeamercolor{itemize item}{fg=CVQCTeal}
  \begin{tikzpicture}
    \path[use as bounding box] (-6.6,-3.15) rectangle (6.6,3.05);

    \begin{scope}[shift={(-5.6,2.85)},x=0.8cm,y=0.47cm]
      % Topology copied from CVQC-paper/figures/brickwork_figure.tex:
% four logical rows and thirteen columns. The caller sets coordinate units.
\begin{scope}[
    qubit/.style={circle,draw=CVQCInk,line width=0.5pt,fill=white,
      inner sep=0pt,minimum size=3.2mm},
    output qubit/.style={qubit,draw=CVQCTeal,fill=CVQCTeal},
    edge/.style={draw=CVQCInk,line width=0.6pt},
    brick/.style={rounded corners=3pt,fill=CVQCPale,
      draw=CVQCMuted!35,line width=0.4pt},
    graph label/.style={font=\footnotesize,text=CVQCMuted}
  ]
  % Each shaded brick spans two rows and five columns.
  \foreach \xa/\xb/\r in {0.5/5.45/1,0.5/5.45/3,4.55/9.45/2,
      8.55/13.5/1,8.55/13.5/3} {
    \draw[brick] (\xa,{-\r+0.38}) rectangle (\xb,{-\r-1.38});
  }

  \foreach \r in {1,...,4} {
    \draw[edge] (1,-\r) -- (13,-\r);
  }
  \foreach \x in {3,5,11,13} {
    \foreach \r in {1,3} {
      \draw[edge] (\x,-\r) -- (\x,{-\r-1});
    }
  }
  \foreach \x in {7,9} {
    \draw[edge] (\x,-2) -- (\x,-3);
  }

  \foreach \r in {1,...,4} {
    \foreach \x in {1,...,12} {
      \node[qubit] at (\x,-\r) {};
    }
    \node[output qubit] at (13,-\r) {};
  }

  \draw[draw=CVQCMuted,line width=0.5pt]
    (0.5,-0.52) -- (0.5,-0.42) -- (5.45,-0.42) -- (5.45,-0.52);
  \node[graph label,anchor=south] at (2.975,-0.42) {one brick};
  \node[graph label,anchor=south] at (13,-0.5) {output};
\end{scope}

    \end{scope}

    \node[anchor=north west,inner sep=0pt,text=CVQCInk] at (-5.35,0.30) {%
      \begin{minipage}{11.2cm}
        \raggedright
        \begin{itemize}
          \setlength{\itemsep}{0.65em}
          \item<1-> Alternative to the circuit model.
          \item<2-> Start with a highly entangled state.
          \item<3-> Compute by adaptive single-qubit measurements.
          \item<4-> Effective operations: $HT^c$ and $\mathsf{CZ}$.
        \end{itemize}
      \end{minipage}%
    };
  \end{tikzpicture}

  \note{
    \footnotesize
    MBQC means measurement-based quantum computation. Show the paper's
    four-row, thirteen-column brickwork graph. Vertices are physical
    qubits and edges represent CZ entangling operations. White vertices
    are measured column by column from left to right; the teal vertices
    in the final column remain as outputs. Shaded rectangles mark the
    repeated two-row, five-column bricks. Each row carries one logical
    qubit through the computation. The prover forms this highly
    entangled graph state from individual qubits; the verifier does not
    have to prepare an entangled state.
    The computation proceeds through single-qubit measurements. Earlier
    outcomes determine later measurement bases and the corrections tracked
    by the verifier. Our effective elementary operations
    are $HT^c$, for $c\in\{0,\ldots,7\}$, and CZ. This family is universal:
    c=0 gives H, composing H with HT gives T, and CZ is entangling.
    MBQC alone does not establish blindness or verification; the encrypted
    construction and tests come later. We will return to the plaintext
    gate-teleportation example in the technical part of the talk.
    Figure source: the paper's Preliminaries, brickwork graph.
  }
\end{frame}
